Sí se simula, y evoluciona. Antes estaba \emph{congelada} en su valor inicial durante los \SI{720}{\second}, mientras el sólido perdía el \qty{28}{\percent} de su masa; de ella dependen la permeabilidad de Kozeny--Carman, la conductividad efectiva del lecho y la corrección de tortuosidad de las difusividades, que quedaban congeladas con ella.

Ahora se recalcula del volumen de sólido que queda, sumando $c_i M_i/\rho_i$ sobre las fases del inventario. Se aplica el cambio \emph{relativo} sobre la porosidad inicial calibrada del caso, no el valor absoluto: la fracción sólida que dan los volúmenes molares (\num{0.413}) no coincide con la que declara el caso (\num{0.46}, de las densidades aparentes), y con la razón esa discrepancia se cancela.

Los valores: parte de \textbf{0.540}, sube a \textbf{0.552} a los \SI{60}{\second} al marcharse los volátiles, y se estabiliza en \textbf{0.678}. Es decir, el lecho pasa de 46.0\,\% de sólido a 32.2\,\%.

El hinchamiento \emph{no} entra en la porosidad: sobre malla fija el lecho no puede crecer de volumen, y lo que el inventario describe es el hueco que dejan los volátiles al marcharse. La expansión se representa aparte.
